\documentclass[
11pt,notheorems,hyperref={pdfauthor=Autor}
]{beamer}
\usepackage{media9}
\addmediapath{figuras}
\addmediapath{figuras/introdução}
\addmediapath{figuras/fundamentação}
\addmediapath{figuras/conclusão}
\addmediapath{figuras/metodologia}
\addmediapath{figuras/resultados}
\input{loadslides.tex}

% =========================
% Metadados da apresentacao
% =========================
% Guia rapido de customizacao do modelo Beamer:
% 1) Ajuste metadados: \title, \subtitle (opcional), \author, \institute, \date.
% 2) Defina orientador/coorientador em \orientador.
% 3) Configure logos da capa em \SetupCustomTitle{instituto}{laboratorio}{universidade}.
% 4) Organize imagens em figuras/ e use \FigureWithFallbackCaption{arquivo}{legenda}.
% 5) Use \pause, \BlockPause e \EquationPause para controlar revelacao de conteudo.
% 6) Estruture a apresentacao com \section e \begin{frame} ... \end{frame}.
\title{Título do trabalho}
\subtitle{Subtítulo do trabalho}
\author{Nome do Autor}
\institute{UNIVERSIDADE ESTADUAL DE CAMPINAS}
\date{Campinas \the\year}

\newcommand{\orientador}{Título e nome do Orientador}

% Capa: Instituto, Laboratório, Universidade
\SetupCustomTitle{logo_feec.png}{logo_dspcom.pdf}{unicamp-logotipo.pdf}

\begin{document}

\begin{frame}[plain]
    \maketitle
\end{frame}

\begin{frame}[plain]{Sumário}
    \tableofcontents
\end{frame}

\section{Introdução}
\begin{frame}{Introdução}

\begin{columns}[T,onlytextwidth]
    \column{0.52\textwidth}
    Coluna 1 

    \pause
    \column{0.48\textwidth}
    Coluna 2
\end{columns}
\end{frame}

\begin{frame}{Introdução}{Exemplo de título de slide}
    Exemplo de imagem

    \centering
    \begin{minipage}{0.8\textwidth}
        \FigureWithFallbackCaption
        {figura_inexistente.png}
        {Exemplo 1.}
    \end{minipage}
\end{frame}

\begin{frame}{Objetivos}
    \pause
    \textbf{Geral}: O objetivo geral da pesquisa.

    \pause
    \textbf{Específicos}:
\begin{itemize}
    \item Os diversos subobjetivos da pesquisa.
\end{itemize}
\end{frame}

\section{Fundamentação}
\begin{frame}{Fundamentação}
\pause
    \begin{itemize}
        \item \textbf{Tópico}: \pause
            \begin{itemize}
                \item Subtópico 1 \pause
                \item Subtópico 2
            \end{itemize}     
    \end{itemize}
\end{frame}


\begin{frame}{Fundamentação}{Aprendizado por Reforço}
Exemplo de coluna dupla, equações e imagem

\begin{columns}[T,onlytextwidth]
    \column{0.50\textwidth}
    
    \Large{$$e^{\pi i}-1 =0$$}
    \LARGE{$$e^{\pi i}-1=0$$}
    \huge{$$e^{\pi i}-1=0$$}
    \Huge{$$e^{\pi i}-1=0$$}
    

    \column{0.50\textwidth}
    \FigureWithFallbackCaption
    {figura_inexiste.png}
    {Exemplo 2.}
\end{columns}
\end{frame}

\section{Metodologia}
\begin{frame}{Metodologia}
Exemplo de coluna dupla de imagens

\pause
    \begin{columns}[T]
    \column{0.25\textwidth}
    \FigureWithFallbackCaption
    {python.png}
    {}
    \pause

    \column{0.25\textwidth}
    \FigureWithFallbackCaption
    {github.png}
    {}
\end{columns}
\end{frame}

\section{Resultados}
\begin{frame}{Resultados}{Experimento II}
Exemplo de coluna dupla de imagens

\begin{columns}[T]
    \pause
    \column{0.45\textwidth}
    \FigureWithFallbackCaption
    {figura_inexistente.pdf}
    {Legenda}

    \pause
    \column{0.45\textwidth}
    \FigureWithFallbackCaption
    {figura_inexistente.pdf}
    {Legenda}
\end{columns}
\end{frame}


\section{Conclusão}
\begin{frame}
    \pause
    \begin{itemize}
        \item Conclusões encontradas com a pesquisa.
    \end{itemize}
\end{frame}


\section{Referências}
\begin{frame}[allowframebreaks]{Referências}
    \small
    \setlength{\leftmargini}{1.8em}
    \setlength{\bibitemsep}{0.25\baselineskip}
    \setlength{\bibhang}{1.0em}
    \raggedright
    \justifying
    \printbibliography
\end{frame}

\end{document}
